% Please do not change %
\documentclass[11pt]{article} % font size
\usepackage[margin=1in]{geometry} % margins
\usepackage{amsmath,amsthm,amssymb} % for the  common "math symbols"
\usepackage[shortlabels]{enumitem} % for enumeration
\usepackage{booktabs} % if you need tables
\usepackage{listings}
\usepackage{color}
\usepackage{tikzit}
\input{sample.tikzstyles}
\input{dots.tikzdefs}

\theoremstyle{definition}
\newtheorem{theorem}{Theorem}[section]
\newtheorem{corollary}[theorem]{Corollary}
\newtheorem{lemma}[theorem]{Lemma}
\newtheorem{proposition}[theorem]{Proposition}
\newtheorem{conjecture}[theorem]{Conjecture}
\newtheorem{definition}[theorem]{Definition}
\newtheorem{fact}[theorem]{Fact}


% Feel free to include additional packages (if needed), but make sure it does not override the margin/font requirements.
%
\setlength{\parindent}{0pt} % no indent for new paragraphs.
\setlength{\parskip}{5pt} % distance between paragraphs, which will be used when you have a blank line in your LaTeX code.


%%%%%%%%%%%%%%%%%%%%%%%%%%%%%%%%%%%%%%%%%%%%%%%%%%%%%%%%%%%%%%%%%%%%%%%%%%%%%%
\title{Quantum Processes and Computation Mini-Project}
%
\author{1633814}
\date{Michaelmas Term 2023}
%
\begin{document}
%
\maketitle

%You may use the usual \LaTeX\ environments to structure your answers:
%\begin{align}
%        \label{basegnn}
%        \mathbf{h}_u^{(t+1)} = \sigma \bigg( \mathbf{W}_\text{self}^{(t)} \mathbf{h}_u^{(t)}  + \mathbf{W}_\text{neigh}^{(t)} \sum_{v \in N(u)} \mathbf{h}_v^{(t)} + \mathbf{b}^{(t)} \bigg). 
%\end{align}


%\begin{table}[h]
%\caption{A simple table.}
%\centering
%\begin{tabular}[t]{lrrrr}
%\toprule
%Graphs & Can & Be  & Fun & Too \\ 
%\midrule 
%$G_1$ & $1$  & $2$ & $3$ & $4$ \\ 
%$G_2$ & $-1$  & $-2$ & $-3$ & $-4$ \\ 
%$G_3$ & $0.1$  & $0.2$ & $0.3$ & $0.4$ \\ 
%\bottomrule
%\end{tabular}
%\label{tab}
%\end{table}

\section*{Question 1}
\addtocounter{section}{1}
The labelled $ZX$-rules shown below \cite{Kissinger_2019} for $\alpha,\beta\in[0,2\pi)$ and $a\in\{0,1\}$ 
are as taught in lectures \cite{coecke2017picturing}, and will be used to justify the equations below.
\setlength{\tabcolsep}{10pt}
\begin{figure}[h!]
\ctikzfig{ZX-rules}
\caption{$ZX$-rules presented in \cite{Kissinger_2019}}\label{table:rewrite}
\end{figure}

% To answer each part of the question, please simply use a paragraph command 
\clearpage
\paragraph{(a)}
% \begin{equation}
% \psi_S = 
% \tikzfig{1a}
% \end{equation}
Let 
\[\psi_S= \tikzfig{psiS}\]
Note for $j,k\in\{0,1\}$, 
\[|j-k|\pi + j\pi + (-1)^k \frac{\pi}{2} = 2n\pi + \frac{\pi}{2}\]
for some $n\in\{0,1\}$.

Then, using $\psi_S$, a 2-qubit measurement, and clasically controlled Pauli $X/Z$ gates, 
the $S$-gate is obtained as desired.
\[\tikzfig{1a}\stackrel{\IdentityRule\PiRule}{=}\tikzfig{1aii}
\stackrel{\SpiderRule\IdentityRule}{=}\tikzfig{1aiii}\stackrel{\SpiderRule\IdentityRule}{=}\tikzfig{1aiv}\]


\paragraph{(b)}
Let
\[\psi_H= \tikzfig{psiH}\]
Then, also using the $S$-gate from part (a),
\[\tikzfig{1b}=\tikzfig{1bii}\stackrel{\HadamardRule}{=}\tikzfig{1biii}
\stackrel{\SpiderRule\IdentityRule\HadamardRule}{=}\tikzfig{1biv}\stackrel{\SpiderRule\IdentityRule}{=}\tikzfig{1bv}\]

\clearpage
\paragraph{(c)}
Let
\[\psi_{CNOT}= \tikzfig{psiC}\]
so we have, using 2 two-qubit measurements, and clasically controlled Pauli $X/Z$ gates,
\begin{equation}
\begin{aligned}
&\tikzfig{1c}\stackrel{\IdentityRule}{=}\tikzfig{1cii}\stackrel{\SpiderRule}{=}\tikzfig{1ciii}\\
&\stackrel{\PiRule}{=}\tikzfig{1civ}\stackrel{\PiRule}{=}\tikzfig{1cv}
\stackrel{\SpiderRule\IdentityRule}{=}\tikzfig{1cvi}
\end{aligned}
\end{equation}

\paragraph{(d)}
No, since the Bell model only has $\circ$ $\frac{n\pi}{2}$-phase gates for some $n\in\{0,1,2,3\}$,
$\circ$ $\frac{\pi}{4}$-phase gates cannot be constructed.

First, we introduce some definitions from \cite{coecke2017picturing}.

\begin{definition}[{\cite{coecke2017picturing} \textbf{Def 9.120.}}][Graph state]
A graph state is a state whose $ZX$-diagram consists only of (0-phase) $\circ$-spiders and $H$-gates where:
\begin{enumerate}
    \item every $\circ$-spider has exactly one output, and
    \item all non-output wires connect two $\circ$-spiders with a single $H$-gate.
\end{enumerate}
\label{def:1}
\end{definition}

\begin{definition}[{\cite{coecke2017picturing} \textbf{Def 9.103.}}][Clifford diagram]
A Clifford diagram is a $ZX$-diagram where the phases are restricted to integer multiples of $\frac{\pi}{2}$.
\label{def:2}
\end{definition}

\begin{definition}[{\cite{coecke2017picturing} \textbf{Def 9.129.}}][Clifford+T diagram]
A Clifford+T diagram is a $ZX$-diagram where the phases are restricted to integer multiples of $\frac{\pi}{4}$.
\label{def:3}
\end{definition}

\begin{definition}[{\cite{coecke2017picturing}}][Local Clifford unitaries]
Local Clifford unitaries are unitaries expressible as parallel compositions of single-system Clifford diagrams.
\label{def:4}
\end{definition}

\begin{proposition}[{\cite{coecke2017picturing} \textbf{Prop 9.122.}}]
Every Clifford diagram can be turned into a Clifford state via process-state duality,
then transformed using $ZX$-calculus into a graph state, followed by local Clifford unitaries.
\label{prop:1}
\end{proposition}

\begin{corollary}
A Clifford diagram is an extension of a graph state which 
can contain integer multiples of $\frac{\pi}{2}$-phase spiders.
\label{cor:1}
\end{corollary}

\begin{corollary}
A Clifford+T diagram is an extension of a Clifford diagram which 
can contain intger multiples of $\frac{\pi}{4}$-phase spiders.
\label{cor:2}
\end{corollary}

\begin{theorem}[{\cite{coecke2017picturing} \textbf{Theorem 9.130.}}]
Clifford+T diagrams are \textit{approximately universal}. 
That is, any linear map can be approximated up to arbitrary precision by a Clifford+T diagram.
\label{thm:1}
\end{theorem}


\begin{definition}[{\cite{coecke2017picturing} \textbf{Section 12.3.1.}}][$CZ$-gate]
The $CZ$-gate can be easily constructed using a CNOT-gate and two Hadamard gates.
\[CZ-gate:=D\quad\tikzfig{1da}\stackrel{\hadamardrule}{=}D\quad\tikzfig{1daii}\]
\end{definition}

\begin{definition}[{\cite{coecke2017picturing} \textbf{Section 12.3.1.}}]
Graph states can be realised by using a simple circuit as follows:
\begin{enumerate}
    \item Prepare $n$ qubits in the $\circ$-state.
    \item To introduce an edge between the $i$-th and $j$-th qubit, apply a $CZ$-gate.
\end{enumerate}
\label{def:5}
\end{definition}

\begin{proposition}
The Bell model is universal for Clifford diagrams.
\end{proposition}
\begin{proof}
Using Proposition \ref{prop:1}, any Clifford diagram 
can be transformed to a graph state with local Clifford unitaries.
Then, from Definition \ref{def:5}, the Bell model has the required CNOT-gates, Hadamard gates, 
and $Z$-basis states (converted to $X$-basis states using the Hadamard) to realise a graph state,
along with the local Clifford unitaries.
\label{prop:2}
\end{proof}

However, Clifford diagrams are not known to be \textit{approximately universal}.
We then require the $T$-gate for Clifford+T diagrams following Corollary \ref{cor:2} to obtain 
\textit{approximate universality} as in Theorem \ref{thm:1}.
\[T:=\tikzfig{pi4}\]

Note for $k\in\{0,1\}$, 
\[\frac{k\pi}{2} + (-1)^k \frac{\pi}{4} = \frac{\pi}{4}\]

Let
\[\psi_T= \tikzfig{psiT}\]
Then we can use a two-qubit measurement, and classically controlled $S$ and Pauli $X/Z$ gates to have
\[\tikzfig{1d}\stackrel{\IdentityRule\PiRule}{=}\tikzfig{1dii}
\stackrel{\SpiderRule\IdentityRule}{=}\tikzfig{1diii}\stackrel{\SpiderRule\IdentityRule}{=}\tikzfig{1div}\]

\paragraph{(e)}

It is just as powerful as before.
Simply use a two-qubit measurement on the $Z$-basis state 
as follows to obtain the equivalent $Z$-basis measurement.
The classical value of the first qubit $(j)$ is essentially discarded.

\[\tikzfig{1e}\stackrel{\copyrule\SpiderRule}{=}\tikzfig{1ei}\]

\clearpage
\section*{Question 2}

\clearpage
\bibliographystyle{plain}
\bibliography{refs}

\end{document}